\section{Related Work}
\label{sec:related}

Chiyoda builds on physical pedestrian-evacuation models rather than replacing
them. Social-force, cellular-automata, and microscopic interaction models
explain how local movement rules can produce bottlenecks, counterflow, density
waves, and other crowd-level failure modes~\cite{helbing1995social,
burstedde2001simulation,moussaid2011simple}. Evacuation-model reviews caution
that such models remain engineering abstractions whose outputs depend on
scenario scope, behavioral assumptions, calibration, and validation
data~\cite{gwynne1999review,schadschneider2009evacuation}. Chiyoda uses this
tradition as the motion substrate; its distinction is the information-control
layer placed on top of that substrate.

Hazard-coupled evacuation and guidance models are closest to the CBRN-like
setting studied here. Toxic-gas and gaseous-material simulations show that
contaminant spread, perception, and information transmission can alter route
choice and evacuation outcomes~\cite{makmul2020cellular,liu2021toxicgas}.
Guidance models similarly show that recommendations must account for hazards,
congestion, and compliance~\cite{chu2017variable}. Chiyoda differs by making
message policy the controlled variable: timing, targeting, budget, credibility,
and reach are varied directly instead of being treated as fixed scenario
assumptions.

Incomplete-information and social-behavior models motivate the belief layer.
Bayesian and game-theoretic route-choice work studies uncertainty and
congestion~\cite{hoogendoorn2004routechoice,wang2023bayesian}, while
social-group ABM surveys and evacuation studies emphasize communication,
affiliation, familiar routes, visible crowd movement, and conflicting social
cues~\cite{templeton2023social,senanayake2024agent,sime1983affiliative,
kobes2010building,bode2013human,kinateder2014social}. Chiyoda keeps these
effects stylized, but exports belief entropy, belief accuracy, reach, exposure,
queueing, and exit imbalance together so belief improvement can be compared
against downstream safety effects.

Existing tools cover much of the surrounding simulation and validation stack.
JuPedSim, Vadere, and SUMO provide mature pedestrian simulation capabilities,
and PedPy provides trajectory-analysis measures for density, velocity, flow,
and fundamental diagrams~\cite{kemlohwagoum2015jupedsim,kleinmeier2019vadere,
alvarezlopez2018sumo,schrodter2025pedpy}. Geometry and hazard studies can draw
on OpenStationMap, GTFS Pathways, and high-fidelity tools such as the NIST Fire
Dynamics Simulator~\cite{openstationmap2026,gtfs2026pathways,mcgrattan2013fds}.
Chiyoda is therefore best read as a replayable information-intervention
benchmark that can be calibrated or cross-checked against these systems, not as
a substitute for full pedestrian or dispersion validation~\cite{ronchi2013verification}.
